\section{AHSME 1961}
        \begin{enumerate}
        	\item (\textit{AHSME 1961}) When simplified $(-\frac{1}{125})^{-2/3}$ becomes:


 $\textbf{(A)}\ \frac{1}{25} \qquad \textbf{(B)}\ -\frac{1}{25} \qquad \textbf{(C)}\ 25\qquad \textbf{(D)}\ -25\qquad \textbf{(E)}\ 25\sqrt{-1}$


	\item (\textit{AHSME 1961}) An automobile travels $a/6$ feet in $r$ seconds. If this rate is maintained for $3$ minutes, how many yards does it travel in $3$ minutes?


 $\textbf{(A)}\ \frac{a}{1080r}\qquad \textbf{(B)}\ \frac{30r}{a}\qquad \textbf{(C)}\ \frac{30a}{r}\qquad \textbf{(D)}\ \frac{10r}{a}\qquad \textbf{(E)}\ \frac{10a}{r}$


	\item (\textit{AHSME 1961}) If the graphs of $2y+x+3=0$ and $3y+ax+2=0$ are to meet at right angles, the value of $a$ is:


 $\textbf{(A)}\ \pm \frac{2}{3} \qquad \textbf{(B)}\ -\frac{2}{3}\qquad \textbf{(C)}\ -\frac{3}{2} \qquad \textbf{(D)}\ 6\qquad \textbf{(E)}\ -6$


	\item (\textit{AHSME 1961}) Let the set consisting of the squares of the positive integers be called $u$; thus $u$ is the set $1, 4, 9, 16 \ldots$. 
If a certain operation on one or more members of the set always yields a member of the set, 
we say that the set is closed under that operation. Then $u$ is closed under:


 $\textbf{(A)}\ \text{Addition}\qquad \textbf{(B)}\ \text{Multiplication} \qquad \textbf{(C)}\ \text{Division} \qquad\\ \textbf{(D)}\ \text{Extraction of a positive integral root}\qquad \textbf{(E)}\ \text{None of these}$


	\item (\textit{AHSME 1961}) Let $S=(x-1)^4+4(x-1)^3+6(x-1)^2+4(x-1)+1$. Then $S$ equals:


 $\textbf{(A)}\ (x-2)^4 \qquad \textbf{(B)}\ (x-1)^4 \qquad \textbf{(C)}\ x^4 \qquad \textbf{(D)}\ (x+1)^4 \qquad \textbf{(E)}\ x^4+1$


	\item (\textit{AHSME 1961}) When simplified, $\log{8} \div \log{\frac{1}{8}}$ becomes:


 $\textbf{(A)}\ 6\log{2} \qquad \textbf{(B)}\ \log{2} \qquad \textbf{(C)}\ 1 \qquad \textbf{(D)}\ 0\qquad \textbf{(E)}\ -1$


	\item (\textit{AHSME 1961}) When simplified, the third term in the expansion of $(\frac{a}{\sqrt{x}}-\frac{\sqrt{x}}{a^2})^6$ is:


 $\textbf{(A)}\ \frac{15}{x}\qquad \textbf{(B)}\ -\frac{15}{x}\qquad \textbf{(C)}\ -\frac{6x^2}{a^9} \qquad \textbf{(D)}\ \frac{20}{a^3}\qquad \textbf{(E)}\ -\frac{20}{a^3}$


	\item (\textit{AHSME 1961}) Let the two base angles of a triangle be $A$ and $B$, with $B$ larger than $A$. 
The altitude to the base divides the vertex angle $C$ into two parts, $C_1$ and $C_2$, with $C_2$ adjacent to side $a$. Then:


 $\textbf{(A)}\ C_1+C_2=A+B \qquad \textbf{(B)}\ C_1-C_2=B-A \qquad\\ \textbf{(C)}\ C_1-C_2=A-B \qquad \textbf{(D)}\ C_1+C_2=B-A\qquad \textbf{(E)}\ C_1-C_2=A+B$


	\item (\textit{AHSME 1961}) Let $r$ be the result of doubling both the base and exponent of $a^b$, and $b$ does not equal to $0$. 
If $r$ equals the product of $a^b$ by $x^b$, then $x$ equals:


 $\textbf{(A)}\ a \qquad \textbf{(B)}\ 2a \qquad \textbf{(C)}\ 4a \qquad \textbf{(D)}\ 2\qquad \textbf{(E)}\ 4$


	\item (\textit{AHSME 1961}) Each side of $\triangle ABC$ is $12$ units. $D$ is the foot of the perpendicular dropped from $A$ on $BC$, 
and $E$ is the midpoint of $AD$. The length of $BE$, in the same unit, is:


 $\textbf{(A)}\ \sqrt{18} \qquad \textbf{(B)}\ \sqrt{28} \qquad \textbf{(C)}\ 6 \qquad \textbf{(D)}\ \sqrt{63} \qquad \textbf{(E)}\ \sqrt{98}$


	\item (\textit{AHSME 1961}) Two tangents are drawn to a circle from an exterior point $A$; they touch the circle at points $B$ and $C$ respectively. 
A third tangent intersects segment $AB$ in $P$ and $AC$ in $R$, and touches the circle at $Q$. If $AB=20$, then the perimeter of triangle $APR$ is:


 $\textbf{(A)}\ 42\qquad \textbf{(B)}\ 40.5 \qquad \textbf{(C)}\ 40\qquad \textbf{(D)}\ 39\frac{7}{8} \qquad \textbf{(E)}\ \text{not determined by the given information}$


	\item (\textit{AHSME 1961}) The first three terms of a geometric progression are $\sqrt{2}, \sqrt[3]{2}, \sqrt[6]{2}$. Find the fourth term.


 $\textbf{(A)}\ 1\qquad \textbf{(B)}\ \sqrt[7]{2}\qquad \textbf{(C)}\ \sqrt[8]{2}\qquad \textbf{(D)}\ \sqrt[9]{2}\qquad \textbf{(E)}\ \sqrt[10]{2}$


	\item (\textit{AHSME 1961}) The symbol $|a|$ means $a$ if $a$ is a positive number or zero, and $-a$ if $a$ is a negative number. 
For all real values of $t$ the expression $\sqrt{t^4+t^2}$ is equal to:


 $\textbf{(A)}\ t^3\qquad \textbf{(B)}\ t^2+t\qquad \textbf{(C)}\ |t^2+t|\qquad \textbf{(D)}\ t\sqrt{t^2+1}\qquad \textbf{(E)}\ |t|\sqrt{1+t^2}$


	\item (\textit{AHSME 1961}) A rhombus is given with one diagonal twice the length of the other diagonal. 
Express the side of the rhombus in terms of $K$, where $K$ is the area of the rhombus in square inches.


 $\textbf{(A)}\ \sqrt{K}\qquad \textbf{(B)}\ \frac{1}{2}\sqrt{2K}\qquad \textbf{(C)}\ \frac{1}{3}\sqrt{3K}\qquad \textbf{(D)}\ \frac{1}{4}\sqrt{4K}\qquad \textbf{(E)}\ \text{None of these are correct}$


	\item (\textit{AHSME 1961}) If $x$ men working $x$ hours a day for each of $x$ days produce $x$ articles, then the number of articles 
(not necessarily an integer) produced by $y$ men working $y$ hours a day for each of $y$ days is:


 $\textbf{(A)}\ \frac{x^3}{y^2}\qquad \textbf{(B)}\ \frac{y^3}{x^2}\qquad \textbf{(C)}\ \frac{x^2}{y^3}\qquad \textbf{(D)}\ \frac{y^2}{x^3}\qquad \textbf{(E)}\ y$


	\item (\textit{AHSME 1961}) An altitude $h$ of a triangle is increased by a length $m$. How much must be taken from the corresponding base $b$ 
so that the area of the new triangle is one-half that of the original triangle? 


 $\textbf{(A)}\ \frac{bm}{h+m}\qquad \textbf{(B)}\ \frac{bh}{2h+2m}\qquad \textbf{(C)}\ \frac{b(2m+h)}{m+h}\qquad \textbf{(D)}\ \frac{b(m+h)}{2m+h}\qquad \textbf{(E)}\ \frac{b(2m+h)}{2(h+m)}$


	\item (\textit{AHSME 1961}) In the base ten number system the number $526$ means $5 \times 10^2+2 \times 10 + 6$. 
In the Land of Mathesis, however, numbers are written in the base $r$. 
Jones purchases an automobile there for $440$ monetary units (abbreviated m.u). 
He gives the salesman a $1000$ m.u bill, and receives, in change, $340$ m.u. The base $r$ is:


 $\textbf{(A)}\ 2\qquad \textbf{(B)}\ 5\qquad \textbf{(C)}\ 7\qquad \textbf{(D)}\ 8\qquad \textbf{(E)}\ 12$


	\item (\textit{AHSME 1961}) The yearly changes in the population census of a town for four consecutive years are, 
respectively, 25\% increase, 25\% increase, 25\% decrease, 25\% decrease. 
The net change over the four years, to the nearest percent, is:


 $\textbf{(A)}\ -12 \qquad \textbf{(B)}\ -1 \qquad \textbf{(C)}\ 0 \qquad \textbf{(D)}\ 1\qquad \textbf{(E)}\ 12$


	\item (\textit{AHSME 1961}) Consider the graphs of $y=2\log{x}$ and $y=\log{2x}$. We may say that:


 $\textbf{(A)}\ \text{They do not intersect}\qquad \\ \textbf{(B)}\ \text{They intersect at 1 point only}\qquad \\ \textbf{(C)}\ \text{They intersect at 2 points only} \qquad \\ \textbf{(D)}\ \text{They intersect at a finite number of points but greater than 2} \qquad \\ \textbf{(E)}\ \text{They coincide}$


	\item (\textit{AHSME 1961}) The set of points satisfying the pair of inequalities $y>2x$ and $y>4-x$ is contained entirely in quadrants:


 $\textbf{(A)}\ \text{I and II}\qquad \textbf{(B)}\ \text{II and III}\qquad \textbf{(C)}\ \text{I and III}\qquad \textbf{(D)}\ \text{III and IV}\qquad \textbf{(E)}\ \text{I and IV}$


	\item (\textit{AHSME 1961}) Medians $AD$ and $CE$ of $\triangle ABC$ intersect in $M$. The midpoint of $AE$ is $N$. 
Let the area of $\triangle MNE$ be $k$ times the area of $\triangle ABC$. Then $k$ equals:


 $\textbf{(A)}\ \frac{1}{6}\qquad \textbf{(B)}\ \frac{1}{8}\qquad \textbf{(C)}\ \frac{1}{9}\qquad \textbf{(D)}\ \frac{1}{12}\qquad \textbf{(E)}\ \frac{1}{16}$


	\item (\textit{AHSME 1961}) If $3x^3-9x^2+kx-12$ is divisible by $x-3$, then it is also divisible by:


 $\textbf{(A)}\ 3x^2-x+4\qquad \textbf{(B)}\ 3x^2-4\qquad \textbf{(C)}\ 3x^2+4\qquad \textbf{(D)}\ 3x-4 \qquad \textbf{(E)}\ 3x+4$


	\item (\textit{AHSME 1961}) Points $P$ and $Q$ are both in the line segment $AB$ and on the same side of its midpoint. $P$ divides $AB$ in the ratio $2:3$, 
and $Q$ divides $AB$ in the ratio $3:4$. If $PQ=2$, then the length of $AB$ is:


 $\textbf{(A)}\ 60\qquad \textbf{(B)}\ 70\qquad \textbf{(C)}\ 75\qquad \textbf{(D)}\ 80\qquad \textbf{(E)}\ 85$


	\item (\textit{AHSME 1961}) Thirty-one books are arranged from left to right in order of increasing prices. 
The price of each book differs by $\textdollar 2$ from that of each adjacent book. 
For the price of the book at the extreme right a customer can buy the middle book and the adjacent one. Then:


 $\textbf{(A)}\ \text{The adjacent book referred to is at the left of the middle book}\qquad \\ \textbf{(B)}\ \text{The middle book sells for \textdollar 36} \qquad \\ \textbf{(C)}\ \text{The cheapest book sells for \textdollar4} \qquad \\ \textbf{(D)}\ \text{The most expensive book sells for \textdollar64 } \qquad \\ \textbf{(E)}\ \text{None of these is correct }$


	\item (\textit{AHSME 1961}) $\triangle ABC$ is isosceles with base $AC$. Points $P$ and $Q$ are respectively in $CB$ and $AB$ and such that $AC=AP=PQ=QB$. 
The number of degrees in $\angle B$ is:


 $\textbf{(A)}\ 25\frac{5}{7}\qquad \textbf{(B)}\ 26\frac{1}{3}\qquad \textbf{(C)}\ 30\qquad \textbf{(D)}\ 40\qquad \textbf{(E)}\ \text{Not determined by the information given}$


	\item (\textit{AHSME 1961}) For a given arithmetic series the sum of the first $50$ terms is $200$, and the sum of the next $50$ terms is $2700$. 
The first term in the series is:


 $\textbf{(A)}\ -1221 \qquad \textbf{(B)}\ -21.5 \qquad \textbf{(C)}\ -20.5 \qquad \textbf{(D)}\ 3 \qquad \textbf{(E)}\ 3.5$


	\item (\textit{AHSME 1961}) Given two equiangular polygons $P_1$ and $P_2$ with different numbers of sides; 
each angle of $P_1$ is $x$ degrees and each angle of $P_2$ is $kx$ degrees, 
where $k$ is an integer greater than $1$. 
The number of possibilities for the pair $(x, k)$ is:


 $\textbf{(A)}\ \infty\qquad \textbf{(B)}\ \text{finite, but greater than 2}\qquad \textbf{(C)}\ 2\qquad \textbf{(D)}\ 1\qquad \textbf{(E)}\ 0$


	\item (\textit{AHSME 1961}) If $2137^{753}$ is multiplied out, the units' digit in the final product is:


 $\textbf{(A)}\ 1\qquad \textbf{(B)}\ 3\qquad \textbf{(C)}\ 5\qquad \textbf{(D)}\ 7\qquad \textbf{(E)}\ 9$


	\item (\textit{AHSME 1961}) Let the roots of $ax^2+bx+c=0$ be $r$ and $s$. The equation with roots $ar+b$ and $as+b$ is:


 $\textbf{(A)}\ x^2-bx-ac=0\qquad \textbf{(B)}\ x^2-bx+ac=0 \qquad\\ \textbf{(C)}\ x^2+3bx+ca+2b^2=0 \qquad \textbf{(D)}\ x^2+3bx-ca+2b^2=0 \qquad\\ \textbf{(E)}\ x^2+bx(2-a)+a^2c+b^2(a+1)=0$


	\item (\textit{AHSME 1961}) If $\log_{10}2=a$ and $\log_{10}3=b$, then $\log_{5}12$ equals:


 $\textbf{(A)}\ \frac{a+b}{1+a}\qquad \textbf{(B)}\ \frac{2a+b}{1+a}\qquad \textbf{(C)}\ \frac{a+2b}{1+a}\qquad \textbf{(D)}\ \frac{2a+b}{1-a}\qquad \textbf{(E)}\ \frac{a+2b}{1-a}$


	\item (\textit{AHSME 1961}) In $\triangle ABC$ the ratio $AC:CB$ is $3:4$. The bisector of the exterior angle at $C$ intersects $BA$ extended at $P$ 
($A$ is between $P$ and $B$). The ratio $PA:AB$ is:


 $\textbf{(A)}\ 1:3 \qquad \textbf{(B)}\ 3:4 \qquad \textbf{(C)}\ 4:3 \qquad \textbf{(D)}\ 3:1 \qquad \textbf{(E)}\ 7:1$


	\item (\textit{AHSME 1961}) A regular polygon of $n$ sides is inscribed in a circle of radius $R$. The area of the polygon is $3R^2$. Then $n$ equals: 


 $\textbf{(A)}\ 8\qquad \textbf{(B)}\ 10\qquad \textbf{(C)}\ 12\qquad \textbf{(D)}\ 15\qquad \textbf{(E)}\ 18$


	\item (\textit{AHSME 1961}) The number of solutions of $2^{2x}-3^{2y}=55$, in which $x$ and $y$ are integers, is:


 $\textbf{(A)}\ 0\qquad \textbf{(B)}\ 1\qquad \textbf{(C)}\ 2\qquad \textbf{(D)}\ 3\qquad \textbf{(E)}\ \text{More than three, but finite}$


	\item (\textit{AHSME 1961}) Let S be the set of values assumed by the function 
 \[\frac{2x+3}{x+2}\]
when $x$ is any member of the interval $x \ge 0$. If there exists a number $M$ such that no number of the set $S$ is greater than $M$, 
then $M$ is an upper bound of $S$. If there exists a number $m$ such that no number of the set $S$ is less than $m$, 
then $m$ is a lower bound of $S$. We may then say:


 $\textbf{(A)}\ \text{m is in S, but M is not in S}\qquad\\ \textbf{(B)}\ \text{M is in S, but m is not in S}\qquad\\ \textbf{(C)}\ \text{Both m and M are in S}\qquad\\ \textbf{(D)}\ \text{Neither m nor M are in S}\qquad\\ \textbf{(E)}\ \text{M does not exist either in or outside S}$


	\item (\textit{AHSME 1961}) The number $695$ is to be written with a factorial base of numeration, that is, $695=a_1+a_2\times2!+a_3\times3!+ \cdots a_n \times n!$ 
where $a_1, a_2, a_3 \cdots a_n$ are integers such that $0 \le a_k \le k$, and $n!$ means $n(n-1)(n-2)\cdots2 \times 1$. Find $a_4$
$\textbf{(A)}\ 0\qquad \textbf{(B)}\ 1\qquad \textbf{(C)}\ 2\qquad \textbf{(D)}\ 3\qquad \textbf{(E)}\ 4$


	\item (\textit{AHSME 1961}) In $\triangle ABC$ the median from $A$ is given perpendicular to the median from $B$. If $BC=7$ and $AC=6$, find the length of $AB$. 


 $\textbf{(A)}\ 4\qquad \textbf{(B)}\ \sqrt{17} \qquad \textbf{(C)}\ 4.25\qquad \textbf{(D)}\ 2\sqrt{5} \qquad \textbf{(E)}\ 4.5$


	\item (\textit{AHSME 1961}) In racing over a given distance $d$ at uniform speed, $A$ can beat $B$ by $20$ yards, $B$ can beat $C$ by $10$ yards, 
and $A$ can beat $C$ by $28$ yards. Then $d$, in yards, equals:


 $\textbf{(A)}\ \text{Not determined by the given information}\qquad \textbf{(B)}\ 58\qquad \textbf{(C)}\ 100\qquad \textbf{(D)}\ 116\qquad \textbf{(E)}\ 120$


	\item (\textit{AHSME 1961}) $\triangle ABC$ is inscribed in a semicircle of radius $r$ so that its base $AB$ coincides with diameter $AB$. 
Point $C$ does not coincide with either $A$ or $B$. Let $s=AC+BC$. Then, for all permissible positions of $C$:


 $\textbf{(A)}\ s^2\le8r^2\qquad \textbf{(B)}\ s^2=8r^2 \qquad \textbf{(C)}\ s^2 \ge 8r^2 \qquad\\ \textbf{(D)}\ s^2\le4r^2 \qquad \textbf{(E)}\ s^2=4r^2$


	\item (\textit{AHSME 1961}) Any five points are taken inside or on a square with side length $1$. Let $a$ be the smallest possible number with the 
property that it is always possible to select one pair of points from these five such that the distance between them 
is equal to or less than $a$. Then $a$ is:


 $\textbf{(A)}\ \sqrt{3}/3\qquad \textbf{(B)}\ \sqrt{2}/2\qquad \textbf{(C)}\ 2\sqrt{2}/3\qquad \textbf{(D)}\ 1 \qquad \textbf{(E)}\ \sqrt{2}$


	\item (\textit{AHSME 1961}) Find the minimum value of $\sqrt{x^2+y^2}$ if $5x+12y=60$. 


 $\textbf{(A)}\ \frac{60}{13}\qquad \textbf{(B)}\ \frac{13}{5}\qquad \textbf{(C)}\ \frac{13}{12}\qquad \textbf{(D)}\ 1\qquad \textbf{(E)}\ 0$


\end{enumerate}